La noción de \enquote{equivalencia} determina si dos conceptos pueden considerarse iguales dentro de un  contexto específico. Es una relación que se establece cuando dichos conceptos comparten características comunes en uno o varios aspectos relevantes.

En la teoría de lenguajes de programación, uno de los conceptos sobre los cuales puede definirse esta relación es el de programa. Un \enquote{programa} es la representación formal de un proceso dinámico, expresado mediante reglas sintácticas establecidas por un lenguaje de programación \cite{abelson1996sicp}; no se limita a ser una secuencia válida de caracteres, sino que incorpora un conjunto definido de comportamientos y  resultados posibles durante su ejecución \cite{lamport1985axiomatic}. 

Estos comportamientos y resultados emergen de las construcciones fundamentales que conforman la mayoría de los programas: \emph{expresiones} y \emph{comandos}. 


Las \enquote{expresiones} corresponden a construcciones que pueden evaluarse y reducirse a un valor, mientras que los \enquote{comandos} representan instrucciones asociadas a modificaciones sobre el estado del programa, generalmente sin producir un valor como resultado \cite{tozzi2025statementexpression}. 
Muchos lenguajes de programación establecen una distinción explícita entre expresiones y comandos, lo que conduce a criterios de equivalencia semántica distintos según la naturaleza de la construcción analizada.


A partir de estas ideas surge el problema ---generalmente indecidible--- de la \emph{equivalencia de programas}: determinar si dos programas son equivalentes bajo la misma definición de equivalencia \cite{lahiri2018program}.


\subsection{Equivalencia en contexto estático}
La \enquote{equivalencia de tipos} es la relación definida por un sistema de tipos para determinar cuándo dos expresiones de tipo representan el mismo tipo \cite{fiveable_type_equivalence}. Esta equivalencia puede clasificarse en dos formas \cite{louden2020programming}:

\begin{itemize}
\item \textbf{Equivalencia estructural.} 
Dos tipos son equivalentes si poseen la misma estructura; es decir, si fueron construidos utilizando los mismos constructores de tipos y sus componentes internos tienen los mismos tipos organizados de la misma manera.

\item \textbf{Equivalencia por nombre.} Dos tipos son equivalentes solo si provienen de la misma declaración de tipo.
\end{itemize}


\subsection{Equivalencia en contexto dinámico }

La \enquote{equivalencia observacional} o \enquote{equivalencia contextual} establece que dos programas son equivalentes si cualquier ocurrencia del primero en un programa completo puede ser reemplazada por el segundo sin afectar los resultados observables de su ejecución \cite{Pitts2002Operational}; es decir, no existe ningún contexto del programa completo que permita distinguir aquel  intercambio entre programas.


En otras palabras, dos programas son equivalentes si, aun teniendo implementaciones internas distintas, exhiben la misma funcionalidad al producir los mismos comportamientos observables ante las mismas entradas.


Lo anterior requiere definir la equivalencia de manera específica para expresiones y comandos: dos expresiones son equivalentes si se evalúan al mismo resultado en todo estado, mientras que dos comandos son equivalentes si, para cualquier estado inicial, ambos divergen o ambos terminan en el mismo estado final \cite{Pierce:SF2}.